\documentclass[11pt, a4paper, oneside]{article}
\usepackage[ngerman]{babel}
\usepackage{worksheet}

\begin{document}
	\author{L. Bung}
	\title{Mehrfache Nullstellen \hspace{10cm} von Polynomfunktionen}
	\subject{Mathematik}
	\class{FHR1}
	\maketitle
	
	\partnertask{Vom Funktionsterm zum Graphen}
	
	Sie haben die Funktion $f(x) = \frac{1}{2}x(x-3)^2$ gegeben.
	
	\begin{enumerate}[label=\alph*)]
		\item Bestimmen Sie den Grad der Funktion $f$.
		\item Bestimmen Sie die Nullstelle(n) von $f$.
		\item Beschreiben Sie, wie der Funktionsgraph von $f$ global verläuft.
		\item Untersuchen Sie, wie sich der Funktionsgraph von $f$ in der Nähe der Nullstelle(n) verhält.
	\end{enumerate}
	
	\checkered[8cm]
	
	\begin{hint}{Mehrfache Nullstellen}
		Hat eine Funktion $f$ die Form $f(x) = (x-x_1)^k \cdot r(x)$ (wobei $r(x)$ wiederum ein Polynom ist), so nennt man $x_1$ eine $k$-fache Nullstelle.
		\vspace{4cm}
	\end{hint}
	
	\partnertask{Vom Graphen zum Funktionsterm}
	
	Bestimmen Sie den zu den Graphen gehörigen Funktionsterm.
	
	\begin{multicols}{2}
		\begin{plot}[xmin=-3, xmax=4, ymin=-1, ymax=3]
			\addplot[domain=-3:4, samples=100, red, thick]{1/4*(x+2)*(x-2)^2};
			\node at (-2.5,2.5) {a)};
		\end{plot}
		
		\begin{plot}[xmin=-3.5, xmax=3.5, ymin=-1, ymax=3]
			\addplot[domain=-3:4, samples=100, red, thick]{1/8*(x+2)^2*(x-2)^2};
			\node at (-2.5,2.5) {b)};
		\end{plot}
	\end{multicols}
	
	\checkered[9.5cm]
	
	\begin{goallist}
		...die Vielfachheit von Nullstellen bei Linearfaktorschreibweise ablesen. &&&\\
		\hline
		...die Vielfachheit der Nullstellen eines Graphen ablesen. &&&\\
		\hline
		...aus den Nullstellen und dem globalen Verlauf eines Graphen auf den Grad der Funktion schließen. &&&\\
	\end{goallist}
\end{document}